\subsection{Definition of vector bundles}

7.1.1. A vector chart of M is a triple \(t = (U, \varphi, F)\) , where U is an open subset of B, F is a Banach space, and \(\varphi\) is a bijection from \(\pi^{-1}(U)\) onto \(U \times F\) such that \(\pi(\varphi^{-1}(b, h)) = b\) for all \(b \in B\) and all \(h \in F\) . We say that U is the domain of the vector chart \(t\) and that \(t\) is a vector chart of M at \(b \in B\) if \(b \in U\) . For every \(b \in U\) , one denotes by \(t_b\) the bijection from F to \(M_b\) defined by \(t_b(h) = \varphi^{-1}(b, h)\) for \(h \in F\) .

7.1.2. Two vector charts \(t = (\mathrm{U}, \varphi, \mathrm{F})\) and \(t' = (\mathrm{U}', \varphi', \mathrm{F}')\) of \(\mathbf{M}\) are \(\mathbf{C}'\) -compatible (or simply compatible) if there exists a map \(\lambda\) of class \(\mathbf{C}'\) from the manifold \(\mathrm{U} \cap \mathrm{U}'\) into the Banach space \(\mathcal{L}(\mathrm{F}; \mathrm{F}')\) such that:

\[
t _ {b} = t _ {b} ^ {\prime} \circ \lambda (b) \quad \text { for all } b \in \mathbf {U} \cap \mathbf {U} ^ {\prime}.
\]

7.1.3. A set of vector charts of M is said to be a C'-vector atlas (or simply vector atlas) of M if it consists of pairwise C'-compatible vector charts whose domains have union B. Two vector atlases A and B of M are said to be C'-equivalent (or equivalent) if A ∪ B is again a vector atlas of M. This relation is an equivalence relation.

7.1.4. A vector bundle structure of class \(C^{r}\) (basis B) on M is the data of a class of equivalent vector atlases (Sets, Chap. II, § 6, no. 9). A vector chart belonging to a vector atlas of this class is called a vector chart of the vector bundle M.

Let M be a vector bundle with base B. For every \(b \in B\) there exists on the fiber \(M_{b}\) a Banach space structure, and only one such that, for every vector chart \(t = (\mathbf{U}, \varphi, \mathbf{F})\) of the vector bundle M at b, the map \(t_{b}\) is an isomorphism from F onto \(M_{b}\) .

Let \(c = (\mathbf{U}, \psi, \mathbf{E})\) be a chart of the manifold \(\mathbf{B}\) and let \(t = (\mathbf{U}, \varphi, \mathbf{F})\) be a vector chart of the vector bundle \(\mathbf{M}\) , with the same domain \(\mathbf{U}\) . For \(x \in \pi^{-1}(\mathbf{U})\) , set:

\[
\alpha (x) = (\psi (\pi (x)), t _ {\pi (x)} ^ {- 1} (x)).
\]

Then, the triple \((\pi^{-1}(\mathbf{U}), \alpha, \mathbf{E} \times \mathbf{F})\) is a chart of the set M. There exists on M a manifold structure of class \(C^{r}\) and only one (called the underlying structure of M) for which all the charts thus obtained are charts of the manifold M. The triple \((\mathbf{M}, \mathbf{B}, \pi)\) is then a fibration (6.1.1).

7.1.5. Let F be a Banach space. Set \(M = B \times F\) , the map being the first projection. There exists on M a vector bundle structure (with base B), and only one, for which \((\mathbf{B}, \mathbf{Id}_{\mathbf{M}}, \mathbf{F})\) is a vector chart. We say that \(B \times F\) endowed with this structure is the trivial vector bundle over the base B with fiber F, and it is sometimes denoted by \(F_{B}\) . The manifold structure of \(F_{B}\) is the product manifold structure, and for every \(b \in B\) , the map \(h \mapsto (b, h)\) is a Banach space isomorphism from F onto the fiber of \(F_{B}\) at the point \(b \in B\) . One often denotes by 0 a trivial vector bundle whose fiber is reduced to 0.

7.1.6. Let M be a vector bundle with base B. For \(b \in B\) we call the dimension (finite or \(+\infty\) ) of the Banach space \(M_{b}\) the rank of M at b, and denote it by \(\operatorname{rg}_{b}(\mathbf{M})\) . One has \(\dim_{x}\mathbf{M} = \dim_{b}\mathbf{B} + \operatorname{rg}_{b}\mathbf{M}\) for \(b = \pi(x)\) . The function \(b \mapsto \operatorname{rg}_{b}\mathbf{M}\) is locally constant. We say that M is of finite rank if \(\operatorname{rg}_{b}\mathbf{M} < +\infty\) for all \(b \in B\) .

